\section{The \texttt{SANS\_Guinier} McStas Component}
Sample for Small Angle Neutron Scattering, Guinier model

\subsection*{Identification}
\begin{itemize}
  \item \textbf{Author:} Henrich Frielinghaus
  \item \textbf{Origin:} FZ-Juelich/FRJ-2/IFF/KWS-2
  \item \textbf{Date:} Sept 2004
\end{itemize}

\subsection*{Description}
Sample that scatters with a Guinier shape. This is just an example where analytically an integral exists. The neutron paths are proportional to the intensity (low intensity \textgreater{} few paths).

Guinier function (Rg) a = Rg*Rg/3 propability\_unscaled = q * exp(-a*q*q) integral\_prop\_unscal = 1/(2*a) * (1 - exp(-a*q*q))

\begin{verbatim}
propability_scaled   = 2*a * q*exp(-a*q*q) / (1 - exp(-a*q*q))
\end{verbatim}

integral\_prop\_scaled = (1 - exp(-a*q*q)) / (1 - exp(-a*qmax*qmax))

In this simulation method many paths occur for high propability. For simulation of low intensities see SANS\_AnySamp.

Sample components leave the units of flux for the probability of the individual paths. That is more consitent than the Sans\_spheres routine. Furthermore one can simulate the transmitted beam. This allows to determine the needed size of the beam stop. Only absorption has not been included yet to these sample-components. But thats really nothing.

Example: SANS\_Guinier(transm=0.5, Rg=100, qmax=0.03, xwidth=0.01, yheight=0.01, zdepth=0.001)

\subsection*{Input parameters}
Parameters in \textbf{boldface} are required; the others are optional.

\begin{longtable}{p{0.22\textwidth}p{0.12\textwidth}p{0.46\textwidth}p{0.14\textwidth}}
\toprule
\textbf{Name} & \textbf{Unit} & \textbf{Description} & \textbf{Default} \\
\midrule
\endhead
transm & 1 & (coherent) transmission of sample for the optical path "zdepth" & 0.5 \\
Rg & \AA{} & Radius of Gyration & 100 \\
qmax & \AA{}$^{-1}$ & Maximum scattering vector & 0.03 \\
xwidth & m & horiz. dimension of sample, as a width & 0.01 \\
yheight & m & vert.. dimension of sample, as a height & 0.01 \\
zdepth & m & thickness of sample & 0.001 \\
\bottomrule
\end{longtable}

\subsection*{Links}
\begin{itemize}
  \item Component source code found in file \texttt{SANS\_Guinier.comp}.
  \item Sans\_spheres component
\end{itemize}
\IfFileExists{contrib/SANS_Guinier_static.tex}{\input{contrib/SANS_Guinier_static.tex}}{}